
\documentclass{article}
\usepackage{colortbl}
\usepackage{makecell}
\usepackage{multirow}
\usepackage{supertabular}

\begin{document}

\newcounter{utterance}

\centering \large Interaction Transcript for game `clean\_up', experiment `2\_hard\_3obj\_el', episode 2 with nemotron{-}3{-}ultra.
\vspace{24pt}

{ \footnotesize  \setcounter{utterance}{1}
\setlength{\tabcolsep}{0pt}
\begin{supertabular}{c@{$\;$}|p{.15\linewidth}@{}p{.15\linewidth}p{.15\linewidth}p{.15\linewidth}p{.15\linewidth}p{.15\linewidth}}
    \# & $\;$A & \multicolumn{4}{c}{Game Master} & $\;\:$B\\
    \hline

    \theutterance \stepcounter{utterance}  
    & & \multicolumn{4}{p{0.6\linewidth}}{
        \cellcolor[rgb]{0.9,0.9,0.9}{
            \makecell[{{p{\linewidth}}}]{
                \texttt{\tiny{[P1$\langle$GM]}}
                \texttt{Είμαι ο αρχηγός του παιχνιδιού σας και παίζετε ένα συνεργατικό παιχνίδι με το ακόλουθο πλέγμα ως πίνακα παιχνιδιού:} \\
\\ 
\\ 
\texttt{ 1234567} \\
\texttt{╔═╤══╤══╗} \\
\texttt{║◌│◌◌│◌◌║ 1} \\
\texttt{║◌└┬─┴──╢ 2} \\
\texttt{║◌◌│L◌◌◌║ 3} \\
\texttt{╟─┬┘◌◌┌─╢ 4} \\
\texttt{║◌│◌◌◌│C║ 5} \\
\texttt{╟─┼─┬─┤P║ 6} \\
\texttt{║◌│◌│◌│◌║ 7} \\
\texttt{╚═╧═╧═╧═╝} \\
\\ 
\\ 
\\ 
\texttt{* Το άνω άκρο εμφανίζει συντεταγμένες x που αυξάνονται προς τα δεξιά και το δεξιό άκρο συντεταγμένες y που αυξάνονται προς τα κάτω.} \\
\texttt{* Τα ακόλουθα αντικείμενα τοποθετούνται τυχαία στο πλέγμα σας: 'C', 'L', 'P'.} \\
\\ 
\\ 
\texttt{Ο άλλος παίκτης βλέπει μια παραλλαγή του σκάφους του παιχνιδιού, όπου τα αντικείμενα τοποθετούνται σε διαφορετικές τυχαίες θέσεις. Δεν μπορείτε να δείτε το ταμπλό του άλλου παίκτη και εκείνος δεν μπορεί να δει το δικό σας.} \\
\\ 
\\ 
\texttt{**Στόχος:**} \\
\\ 
\texttt{Και οι δύο παίκτες πρέπει να μετακινήσουν τα αντικείμενα στο αντίστοιχο φόντο τους, έτσι ώστε πανομοιότυπα αντικείμενα να καταλήξουν στις ίδιες συντεταγμένες. Πρέπει να επικοινωνήσετε με τον άλλο παίκτη για να συμφωνήσετε σε μια κοινή διαμόρφωση του στόχου.} \\
\\ 
\\ 
\texttt{**Κανόνες:**} \\
\\ 
\texttt{* Σε κάθε γύρο, μπορείτε να στείλετε ακριβώς μία από τις ακόλουθες δύο εντολές:} \\
\texttt{ 1. `ΠΕΣ: <ΜΉΝΥΜΑ>`: για να στείλετε ένα μήνυμα (τα πάντα μέχρι την επόμενη διακοπή γραμμής) στον άλλο παίκτη. Θα το προωθήσω στον παρτενέρ σας.} \\
\texttt{2. `ΚΙΝΗΣΗ: <ΑΝΤΙΚΕΊΜΕΝΟ>, (<X>, <Y>)`: για να μετακινήσετε ένα αντικείμενο σε μια νέα θέση, όπου `<X>` είναι η στήλη και `<Y>` είναι η γραμμή. Θα σας ενημερώσω αν η μετακίνησή σας ήταν επιτυχής ή όχι.} \\
\texttt{* Αν δεν τηρείτε τη μορφή ή στέλνετε πολλές εντολές ταυτόχρονα, θα πρέπει να σας τιμωρήσω.} \\
\texttt{* Εάν και οι δύο παίκτες συγκεντρώσουν περισσότερες από 8 ποινές, χάνετε και οι δύο το παιχνίδι.} \\
\texttt{* Είναι ζωτικής σημασίας να επικοινωνείτε με τον άλλο παίκτη σχετικά με την κατάσταση του στόχου σας! Ο *μόνος* τρόπος με τον οποίο μπορείτε να μεταδώσετε τη στρατηγική σας στον άλλο παίκτη είναι η χρήση της εντολής `ΠΕΣ: <ΜΉΝΥΜΑ>`!} \\
\\ 
\\ 
\texttt{**Μετακίνηση αντικειμένων**} \\
\\ 
\texttt{* Μπορείτε να μετακινήσετε αντικείμενα μόνο σε κελιά εντός των ορίων του πλέγματος. Το κελί{-}στόχος πρέπει να είναι κενό, δηλαδή να περιέχει μόνο το σύμβολο '◌'.} \\
\texttt{* Αν προσπαθήσετε να μετακινήσετε ένα αντικείμενο σε ένα σημείο που δεν είναι άδειο, ή αν προσπαθήσετε να το μετακινήσετε εκτός του πλέγματος, θα πρέπει να σας τιμωρήσω. Μπορείτε να προσπαθήσετε ξανά.} \\
\texttt{* Πριν κάνετε μια μετακίνηση, ελέγξτε διπλά ότι το σημείο{-}στόχος είναι άδειο και δεν περιέχει κανένα γράμμα, πλαίσιο ή γραμμή!} \\
\\ 
\\ 
\texttt{**Τέλος του παιχνιδιού**} \\
\\ 
\texttt{Αν νομίζετε ότι έχετε πετύχει το στόχο της ευθυγράμμισης όλων των αντικειμένων, μπορείτε να ζητήσετε από τον άλλο παίκτη να τερματίσει το παιχνίδι στέλνοντας `ΠΕΣ: τελείωσες;`. Αν ο άλλος παίκτης σας ζητήσει να τερματίσετε το παιχνίδι και απαντήσετε με `ΠΕΣ: τελείωσα!`, το παιχνίδι θα τελειώσει.} \\
\\ 
\texttt{Και οι δύο παίκτες κερδίζουν αν το παιχνίδι τελειώσει μέσα σε 12 γύρους, όπου ένας γύρος ορίζεται ως δύο παίκτες που στέλνουν από μία έγκυρη εντολή.} \\
\\ 
\texttt{**Βαθμολογία:**} \\
\\ 
\texttt{Όσο πιο κοντά βρίσκονται τα ίδια αντικείμενα και στους δύο πίνακες του παιχνιδιού, τόσο περισσότερους πόντους παίρνετε. Οι ποινές μειώνουν τους πόντους σας. Μπορείτε να καταρρίψετε το ρεκόρ;} \\
\\ 
\\ 
\texttt{Παρακαλώ στείλτε ένα μήνυμα στον άλλο παίκτη για να ξεκινήσετε το παιχνίδι!} \\
            }
        }
    }
    & \\ \\

    \theutterance \stepcounter{utterance}  
    & \multicolumn{4}{p{0.6\linewidth}}{
        \cellcolor[rgb]{0.9,0.9,0.9}{
            \makecell[{{p{\linewidth}}}]{
                \texttt{\tiny{[P1$\rangle$GM]}}
                \texttt{ΠΕΣ: Γεια! Ας ξεκινήσουμε. Έχω 3 αντικείμενα: C (στην κελιά 7,5), L (στην κελιά 3,3) και P (στην κελιά 4,6). Πού βρίσκονται τα δικά σου;} \\
            }
        }
    }
    & & \\ \\

    \theutterance \stepcounter{utterance}  
    & & & \multicolumn{2}{p{0.3\linewidth}}{
        \cellcolor[rgb]{0.9,0.9,0.9}{
            \makecell[{{p{\linewidth}}}]{
                \texttt{\tiny{[GM$|$GM]}}
                \texttt{ΠΕΣ: Γεια! Ας ξεκινήσουμε. Έχω 3 αντικείμενα: C (στην κελιά 7,5), L (στην κελιά 3,3) και P (στην κελιά 4,6). Πού βρίσκονται τα δικά σου;} \\
            }
        }
    }
    & & \\ \\

    \theutterance \stepcounter{utterance}  
    & & \multicolumn{4}{p{0.6\linewidth}}{
        \cellcolor[rgb]{0.9,0.9,0.9}{
            \makecell[{{p{\linewidth}}}]{
                \texttt{\tiny{[GM$\rangle$P2]}}
                \texttt{Είμαι ο αρχηγός του παιχνιδιού σας και παίζετε ένα συνεργατικό παιχνίδι με το ακόλουθο πλέγμα ως πίνακα παιχνιδιού:} \\
\\ 
\\ 
\texttt{ 1234567} \\
\texttt{╔═╤══╤══╗} \\
\texttt{║◌│◌◌│◌◌║ 1} \\
\texttt{║L└┬─┴──╢ 2} \\
\texttt{║◌C│◌◌◌◌║ 3} \\
\texttt{╟─┬┘◌◌┌─╢ 4} \\
\texttt{║◌│◌◌◌│◌║ 5} \\
\texttt{╟─┼─┬─┤◌║ 6} \\
\texttt{║◌│P│◌│◌║ 7} \\
\texttt{╚═╧═╧═╧═╝} \\
\\ 
\\ 
\\ 
\texttt{* Το άνω άκρο εμφανίζει συντεταγμένες x που αυξάνονται προς τα δεξιά και το δεξιό άκρο συντεταγμένες y που αυξάνονται προς τα κάτω.} \\
\texttt{* Τα ακόλουθα αντικείμενα τοποθετούνται τυχαία στο πλέγμα σας: 'C', 'L', 'P'.} \\
\\ 
\\ 
\texttt{Ο άλλος παίκτης βλέπει μια παραλλαγή του σκάφους του παιχνιδιού, όπου τα αντικείμενα τοποθετούνται σε διαφορετικές τυχαίες θέσεις. Δεν μπορείτε να δείτε το ταμπλό του άλλου παίκτη και εκείνος δεν μπορεί να δει το δικό σας.} \\
\\ 
\\ 
\texttt{**Στόχος:**} \\
\\ 
\texttt{Και οι δύο παίκτες πρέπει να μετακινήσουν τα αντικείμενα στο αντίστοιχο φόντο τους, έτσι ώστε πανομοιότυπα αντικείμενα να καταλήξουν στις ίδιες συντεταγμένες. Πρέπει να επικοινωνήσετε με τον άλλο παίκτη για να συμφωνήσετε σε μια κοινή διαμόρφωση του στόχου.} \\
\\ 
\\ 
\texttt{**Κανόνες:**} \\
\\ 
\texttt{* Σε κάθε γύρο, μπορείτε να στείλετε ακριβώς μία από τις ακόλουθες δύο εντολές:} \\
\texttt{ 1. `ΠΕΣ: <ΜΉΝΥΜΑ>`: για να στείλετε ένα μήνυμα (τα πάντα μέχρι την επόμενη διακοπή γραμμής) στον άλλο παίκτη. Θα το προωθήσω στον παρτενέρ σας.} \\
\texttt{2. `ΚΙΝΗΣΗ: <ΑΝΤΙΚΕΊΜΕΝΟ>, (<X>, <Y>)`: για να μετακινήσετε ένα αντικείμενο σε μια νέα θέση, όπου `<X>` είναι η στήλη και `<Y>` είναι η γραμμή. Θα σας ενημερώσω αν η μετακίνησή σας ήταν επιτυχής ή όχι.} \\
\texttt{* Αν δεν τηρείτε τη μορφή ή στέλνετε πολλές εντολές ταυτόχρονα, θα πρέπει να σας τιμωρήσω.} \\
\texttt{* Εάν και οι δύο παίκτες συγκεντρώσουν περισσότερες από 8 ποινές, χάνετε και οι δύο το παιχνίδι.} \\
\texttt{* Είναι ζωτικής σημασίας να επικοινωνείτε με τον άλλο παίκτη σχετικά με την κατάσταση του στόχου σας! Ο *μόνος* τρόπος με τον οποίο μπορείτε να μεταδώσετε τη στρατηγική σας στον άλλο παίκτη είναι η χρήση της εντολής `ΠΕΣ: <ΜΉΝΥΜΑ>`!} \\
\\ 
\\ 
\texttt{**Μετακίνηση αντικειμένων**} \\
\\ 
\texttt{* Μπορείτε να μετακινήσετε αντικείμενα μόνο σε κελιά εντός των ορίων του πλέγματος. Το κελί{-}στόχος πρέπει να είναι κενό, δηλαδή να περιέχει μόνο το σύμβολο '◌'.} \\
\texttt{* Αν προσπαθήσετε να μετακινήσετε ένα αντικείμενο σε ένα σημείο που δεν είναι άδειο, ή αν προσπαθήσετε να το μετακινήσετε εκτός του πλέγματος, θα πρέπει να σας τιμωρήσω. Μπορείτε να προσπαθήσετε ξανά.} \\
\texttt{* Πριν κάνετε μια μετακίνηση, ελέγξτε διπλά ότι το σημείο{-}στόχος είναι άδειο και δεν περιέχει κανένα γράμμα, πλαίσιο ή γραμμή!} \\
\\ 
\\ 
\texttt{**Τέλος του παιχνιδιού**} \\
\\ 
\texttt{Αν νομίζετε ότι έχετε πετύχει το στόχο της ευθυγράμμισης όλων των αντικειμένων, μπορείτε να ζητήσετε από τον άλλο παίκτη να τερματίσει το παιχνίδι στέλνοντας `ΠΕΣ: τελείωσες;`. Αν ο άλλος παίκτης σας ζητήσει να τερματίσετε το παιχνίδι και απαντήσετε με `ΠΕΣ: τελείωσα!`, το παιχνίδι θα τελειώσει.} \\
\\ 
\texttt{Και οι δύο παίκτες κερδίζουν αν το παιχνίδι τελειώσει μέσα σε 12 γύρους, όπου ένας γύρος ορίζεται ως δύο παίκτες που στέλνουν από μία έγκυρη εντολή.} \\
\\ 
\texttt{**Βαθμολογία:**} \\
\\ 
\texttt{Όσο πιο κοντά βρίσκονται τα ίδια αντικείμενα και στους δύο πίνακες του παιχνιδιού, τόσο περισσότερους πόντους παίρνετε. Οι ποινές μειώνουν τους πόντους σας. Μπορείτε να καταρρίψετε το ρεκόρ;} \\
\\ 
\texttt{Ο άλλος παίκτης ξεκίνησε το παιχνίδι στέλνοντας αυτό το μήνυμα:} \\
\\ 
\texttt{"Γεια! Ας ξεκινήσουμε. Έχω 3 αντικείμενα: C (στην κελιά 7,5), L (στην κελιά 3,3) και P (στην κελιά 4,6). Πού βρίσκονται τα δικά σου;"} \\
\\ 
\texttt{Ποια είναι η πρώτη σας εντολή;} \\
            }
        }
    }
    & \\ \\

    \theutterance \stepcounter{utterance}  
    & & & \multicolumn{4}{p{0.6\linewidth}}{
        \cellcolor[rgb]{0.9,0.9,0.9}{
            \makecell[{{p{\linewidth}}}]{
                \texttt{\tiny{[GM$\langle$P2]}}
                \texttt{ΣΕ: Έχω C στη (5,3), L στη (2,2) και P στη (6,7). Πού θες να τα στοιχιστούμε; Προτείνω να τα βάλουμε στην πρώτη σειρά για ευκολία: C στη (1,1), L στη (2,1), P στη (3,1). Σου φαίνεται καλό;} \\
            }
        }
    }
     \\ \\

    \theutterance \stepcounter{utterance}  
    & & & \multicolumn{2}{p{0.3\linewidth}}{
        \cellcolor[rgb]{0.9,0.9,0.9}{
            \makecell[{{p{\linewidth}}}]{
                \texttt{\tiny{[GM$|$GM]}}
                \texttt{ΣΕ: Έχω C στη (5,3), L στη (2,2) και P στη (6,7). Πού θες να τα στοιχιστούμε; Προτείνω να τα βάλουμε στην πρώτη σειρά για ευκολία: C στη (1,1), L στη (2,1), P στη (3,1). Σου φαίνεται καλό;} \\
            }
        }
    }
    & & \\ \\

    \theutterance \stepcounter{utterance}  
    & & & \multicolumn{2}{p{0.3\linewidth}}{
        \cellcolor[rgb]{0.9,0.9,0.9}{
            \makecell[{{p{\linewidth}}}]{
                \texttt{\tiny{[GM$|$GM]}}
                \texttt{Invalid response format} \\
            }
        }
    }
    & & \\ \\

    \theutterance \stepcounter{utterance}  
    & & \multicolumn{4}{p{0.6\linewidth}}{
        \cellcolor[rgb]{0.9,0.9,0.9}{
            \makecell[{{p{\linewidth}}}]{
                \texttt{\tiny{[GM$\rangle$P2]}}
                \texttt{Ποινή: Το μήνυμά σας δεν είναι στην αναμενόμενη μορφή!} \\
\\ 
\texttt{Βεβαιωθείτε ότι η απάντησή σας περιέχει μόνο `ΠΕΣ: <ΜΉΝΥΜΑ>` ή `ΚΙΝΗΣΗ: <ΑΝΤΙΚΕΊΜΕΝΟ>, (<X>, <Y>)`, και τίποτα άλλο!} \\
\\ 
\texttt{Έχετε συσσωρεύσει συλλογικά 1 από 8 ποινές.} \\
\texttt{Παρακαλώ δοκιμάστε ξανά!} \\
            }
        }
    }
    & \\ \\

    \theutterance \stepcounter{utterance}  
    & & & \multicolumn{4}{p{0.6\linewidth}}{
        \cellcolor[rgb]{0.9,0.9,0.9}{
            \makecell[{{p{\linewidth}}}]{
                \texttt{\tiny{[GM$\langle$P2]}}
                \texttt{ΠΕΣ: Έχω C στη (5,3), L στη (2,2) και P στη (6,7). Πού θες να τα στοιχιστούμε; Προτείνω πρώτη σειρά: C (1,1), L (2,1), P (3,1). Σου φαίνεται καλό;} \\
            }
        }
    }
     \\ \\

    \theutterance \stepcounter{utterance}  
    & & & \multicolumn{2}{p{0.3\linewidth}}{
        \cellcolor[rgb]{0.9,0.9,0.9}{
            \makecell[{{p{\linewidth}}}]{
                \texttt{\tiny{[GM$|$GM]}}
                \texttt{ΠΕΣ: Έχω C στη (5,3), L στη (2,2) και P στη (6,7). Πού θες να τα στοιχιστούμε; Προτείνω πρώτη σειρά: C (1,1), L (2,1), P (3,1). Σου φαίνεται καλό;} \\
            }
        }
    }
    & & \\ \\

    \theutterance \stepcounter{utterance}  
    & & \multicolumn{4}{p{0.6\linewidth}}{
        \cellcolor[rgb]{0.9,0.9,0.9}{
            \makecell[{{p{\linewidth}}}]{
                \texttt{\tiny{[P1$\langle$GM]}}
                \texttt{Το μήνυμά σας μεταφέρθηκε στον άλλο παίκτη.} \\
\\ 
\\ 
\texttt{Αυτή τη στιγμή παίζετε τον γύρο 1 από συνολικά 12.} \\
\texttt{Έχετε συσσωρεύσει συλλογικά 1 από 8 ποινές.} \\
\texttt{Ο άλλος παίκτης έστειλε αυτό το μήνυμα:} \\
\\ 
\\ 
\texttt{"Έχω C στη (5,3), L στη (2,2) και P στη (6,7). Πού θες να τα στοιχιστούμε; Προτείνω πρώτη σειρά: C (1,1), L (2,1), P (3,1). Σου φαίνεται καλό;"} \\
\\ 
\\ 
\texttt{Ποια είναι η επόμενη εντολή σας;} \\
            }
        }
    }
    & \\ \\

    \theutterance \stepcounter{utterance}  
    & \multicolumn{4}{p{0.6\linewidth}}{
        \cellcolor[rgb]{0.9,0.9,0.9}{
            \makecell[{{p{\linewidth}}}]{
                \texttt{\tiny{[P1$\rangle$GM]}}
                \texttt{ΠΕΣ: Εντάξει, ας στοιχιστούμε στην πρώτη σειρά. Ωστόσο, το (1,1) είναι άκυρο (έχει πλαίσιο). Προτείνω τα εξής ελεύθερα κελιά: C {-}> (2,1), L {-}> (3,1), P {-}> (4,1). Σου φαίνεται καλό;} \\
            }
        }
    }
    & & \\ \\

    \theutterance \stepcounter{utterance}  
    & & & \multicolumn{2}{p{0.3\linewidth}}{
        \cellcolor[rgb]{0.9,0.9,0.9}{
            \makecell[{{p{\linewidth}}}]{
                \texttt{\tiny{[GM$|$GM]}}
                \texttt{ΠΕΣ: Εντάξει, ας στοιχιστούμε στην πρώτη σειρά. Ωστόσο, το (1,1) είναι άκυρο (έχει πλαίσιο). Προτείνω τα εξής ελεύθερα κελιά: C {-}> (2,1), L {-}> (3,1), P {-}> (4,1). Σου φαίνεται καλό;} \\
            }
        }
    }
    & & \\ \\

    \theutterance \stepcounter{utterance}  
    & & \multicolumn{4}{p{0.6\linewidth}}{
        \cellcolor[rgb]{0.9,0.9,0.9}{
            \makecell[{{p{\linewidth}}}]{
                \texttt{\tiny{[GM$\rangle$P2]}}
                \texttt{Το μήνυμά σας μεταφέρθηκε στον άλλο παίκτη.} \\
\\ 
\\ 
\texttt{Αυτή τη στιγμή παίζετε τον γύρο 2 από συνολικά 12.} \\
\texttt{Έχετε συσσωρεύσει συλλογικά 1 από 8 ποινές.} \\
\texttt{Ο άλλος παίκτης έστειλε αυτό το μήνυμα:} \\
\\ 
\\ 
\texttt{"Εντάξει, ας στοιχιστούμε στην πρώτη σειρά. Ωστόσο, το (1,1) είναι άκυρο (έχει πλαίσιο). Προτείνω τα εξής ελεύθερα κελιά: C {-}> (2,1), L {-}> (3,1), P {-}> (4,1). Σου φαίνεται καλό;"} \\
\\ 
\\ 
\texttt{Ποια είναι η επόμενη εντολή σας;} \\
            }
        }
    }
    & \\ \\

    \theutterance \stepcounter{utterance}  
    & & & \multicolumn{4}{p{0.6\linewidth}}{
        \cellcolor[rgb]{0.9,0.9,0.9}{
            \makecell[{{p{\linewidth}}}]{
                \texttt{\tiny{[GM$\langle$P2]}}
                \texttt{ΠΕΣ: Συμφωνώ. Στόχοι: C {-}> (2,1), L {-}> (3,1), P {-}> (4,1). Ξεκινάω με το C.} \\
            }
        }
    }
     \\ \\

    \theutterance \stepcounter{utterance}  
    & & & \multicolumn{2}{p{0.3\linewidth}}{
        \cellcolor[rgb]{0.9,0.9,0.9}{
            \makecell[{{p{\linewidth}}}]{
                \texttt{\tiny{[GM$|$GM]}}
                \texttt{ΠΕΣ: Συμφωνώ. Στόχοι: C {-}> (2,1), L {-}> (3,1), P {-}> (4,1). Ξεκινάω με το C.} \\
            }
        }
    }
    & & \\ \\

    \theutterance \stepcounter{utterance}  
    & & \multicolumn{4}{p{0.6\linewidth}}{
        \cellcolor[rgb]{0.9,0.9,0.9}{
            \makecell[{{p{\linewidth}}}]{
                \texttt{\tiny{[P1$\langle$GM]}}
                \texttt{Το μήνυμά σας μεταφέρθηκε στον άλλο παίκτη.} \\
\\ 
\\ 
\texttt{Αυτή τη στιγμή παίζετε τον γύρο 2 από συνολικά 12.} \\
\texttt{Έχετε συσσωρεύσει συλλογικά 1 από 8 ποινές.} \\
\texttt{Ο άλλος παίκτης έστειλε αυτό το μήνυμα:} \\
\\ 
\\ 
\texttt{"Συμφωνώ. Στόχοι: C {-}> (2,1), L {-}> (3,1), P {-}> (4,1). Ξεκινάω με το C."} \\
\\ 
\\ 
\texttt{Ποια είναι η επόμενη εντολή σας;} \\
            }
        }
    }
    & \\ \\

    \theutterance \stepcounter{utterance}  
    & \multicolumn{4}{p{0.6\linewidth}}{
        \cellcolor[rgb]{0.9,0.9,0.9}{
            \makecell[{{p{\linewidth}}}]{
                \texttt{\tiny{[P1$\rangle$GM]}}
                \texttt{ΚΙΝΗΣΗ: C, (2, 1)} \\
            }
        }
    }
    & & \\ \\

    \theutterance \stepcounter{utterance}  
    & & & \multicolumn{2}{p{0.3\linewidth}}{
        \cellcolor[rgb]{0.9,0.9,0.9}{
            \makecell[{{p{\linewidth}}}]{
                \texttt{\tiny{[GM$|$GM]}}
                \texttt{ΚΙΝΗΣΗ: C, (2, 1)} \\
            }
        }
    }
    & & \\ \\

    \theutterance \stepcounter{utterance}  
    & & & \multicolumn{2}{p{0.3\linewidth}}{
        \cellcolor[rgb]{0.9,0.9,0.9}{
            \makecell[{{p{\linewidth}}}]{
                \texttt{\tiny{[GM$|$GM]}}
                \texttt{Ποινή: Το (2,1) δεν είναι κενό, αλλά περιέχει '│'.} \\
\\ 
\texttt{Έχετε συσσωρεύσει συλλογικά 2 από 8 ποινές.} \\
\texttt{Παρακαλώ δοκιμάστε ξανά!} \\
            }
        }
    }
    & & \\ \\

    \theutterance \stepcounter{utterance}  
    & & \multicolumn{4}{p{0.6\linewidth}}{
        \cellcolor[rgb]{0.9,0.9,0.9}{
            \makecell[{{p{\linewidth}}}]{
                \texttt{\tiny{[P1$\langle$GM]}}
                \texttt{Ποινή: Το (2,1) δεν είναι κενό, αλλά περιέχει '│'.} \\
\\ 
\texttt{Έχετε συσσωρεύσει συλλογικά 2 από 8 ποινές.} \\
\texttt{Παρακαλώ δοκιμάστε ξανά!} \\
            }
        }
    }
    & \\ \\

    \theutterance \stepcounter{utterance}  
    & \multicolumn{4}{p{0.6\linewidth}}{
        \cellcolor[rgb]{0.9,0.9,0.9}{
            \makecell[{{p{\linewidth}}}]{
                \texttt{\tiny{[P1$\rangle$GM]}}
                \texttt{ΠΕΣ: Το (2,1) έχει τοίχιο. Ας προσπαθήσουμε στη σειρά 2: C {-}> (2,2), L {-}> (3,2), P {-}> (4,2). Εντάξει;} \\
            }
        }
    }
    & & \\ \\

    \theutterance \stepcounter{utterance}  
    & & & \multicolumn{2}{p{0.3\linewidth}}{
        \cellcolor[rgb]{0.9,0.9,0.9}{
            \makecell[{{p{\linewidth}}}]{
                \texttt{\tiny{[GM$|$GM]}}
                \texttt{ΠΕΣ: Το (2,1) έχει τοίχιο. Ας προσπαθήσουμε στη σειρά 2: C {-}> (2,2), L {-}> (3,2), P {-}> (4,2). Εντάξει;} \\
            }
        }
    }
    & & \\ \\

    \theutterance \stepcounter{utterance}  
    & & \multicolumn{4}{p{0.6\linewidth}}{
        \cellcolor[rgb]{0.9,0.9,0.9}{
            \makecell[{{p{\linewidth}}}]{
                \texttt{\tiny{[GM$\rangle$P2]}}
                \texttt{Το μήνυμά σας μεταφέρθηκε στον άλλο παίκτη.} \\
\\ 
\\ 
\texttt{Αυτή τη στιγμή παίζετε τον γύρο 3 από συνολικά 12.} \\
\texttt{Έχετε συσσωρεύσει συλλογικά 2 από 8 ποινές.} \\
\texttt{Ο άλλος παίκτης έστειλε αυτό το μήνυμα:} \\
\\ 
\\ 
\texttt{"Το (2,1) έχει τοίχιο. Ας προσπαθήσουμε στη σειρά 2: C {-}> (2,2), L {-}> (3,2), P {-}> (4,2). Εντάξει;"} \\
\\ 
\\ 
\texttt{Ποια είναι η επόμενη εντολή σας;} \\
            }
        }
    }
    & \\ \\

    \theutterance \stepcounter{utterance}  
    & & & \multicolumn{4}{p{0.6\linewidth}}{
        \cellcolor[rgb]{0.9,0.9,0.9}{
            \makecell[{{p{\linewidth}}}]{
                \texttt{\tiny{[GM$\langle$P2]}}
                \texttt{ΠΕΣ: Έχεις δίκιο, το (2,1) έχει τοίχιο. Η σειρά 2 είναι ελεύθερη. Συμφωνώ: C {-}> (2,2), L {-}> (3,2), P {-}> (4,2). Ξεκινάω με τη μετακίνηση του C.} \\
            }
        }
    }
     \\ \\

    \theutterance \stepcounter{utterance}  
    & & & \multicolumn{2}{p{0.3\linewidth}}{
        \cellcolor[rgb]{0.9,0.9,0.9}{
            \makecell[{{p{\linewidth}}}]{
                \texttt{\tiny{[GM$|$GM]}}
                \texttt{ΠΕΣ: Έχεις δίκιο, το (2,1) έχει τοίχιο. Η σειρά 2 είναι ελεύθερη. Συμφωνώ: C {-}> (2,2), L {-}> (3,2), P {-}> (4,2). Ξεκινάω με τη μετακίνηση του C.} \\
            }
        }
    }
    & & \\ \\

    \theutterance \stepcounter{utterance}  
    & & \multicolumn{4}{p{0.6\linewidth}}{
        \cellcolor[rgb]{0.9,0.9,0.9}{
            \makecell[{{p{\linewidth}}}]{
                \texttt{\tiny{[P1$\langle$GM]}}
                \texttt{Το μήνυμά σας μεταφέρθηκε στον άλλο παίκτη.} \\
\\ 
\\ 
\texttt{Αυτή τη στιγμή παίζετε τον γύρο 3 από συνολικά 12.} \\
\texttt{Έχετε συσσωρεύσει συλλογικά 2 από 8 ποινές.} \\
\texttt{Ο άλλος παίκτης έστειλε αυτό το μήνυμα:} \\
\\ 
\\ 
\texttt{"Έχεις δίκιο, το (2,1) έχει τοίχιο. Η σειρά 2 είναι ελεύθερη. Συμφωνώ: C {-}> (2,2), L {-}> (3,2), P {-}> (4,2). Ξεκινάω με τη μετακίνηση του C."} \\
\\ 
\\ 
\texttt{Ποια είναι η επόμενη εντολή σας;} \\
            }
        }
    }
    & \\ \\

    \theutterance \stepcounter{utterance}  
    & \multicolumn{4}{p{0.6\linewidth}}{
        \cellcolor[rgb]{0.9,0.9,0.9}{
            \makecell[{{p{\linewidth}}}]{
                \texttt{\tiny{[P1$\rangle$GM]}}
                \texttt{ΚΙΝΗΣΗ: C, (2, 2)} \\
            }
        }
    }
    & & \\ \\

    \theutterance \stepcounter{utterance}  
    & & & \multicolumn{2}{p{0.3\linewidth}}{
        \cellcolor[rgb]{0.9,0.9,0.9}{
            \makecell[{{p{\linewidth}}}]{
                \texttt{\tiny{[GM$|$GM]}}
                \texttt{ΚΙΝΗΣΗ: C, (2, 2)} \\
            }
        }
    }
    & & \\ \\

    \theutterance \stepcounter{utterance}  
    & & & \multicolumn{2}{p{0.3\linewidth}}{
        \cellcolor[rgb]{0.9,0.9,0.9}{
            \makecell[{{p{\linewidth}}}]{
                \texttt{\tiny{[GM$|$GM]}}
                \texttt{Ποινή: Το (2,2) δεν είναι κενό, αλλά περιέχει '└'.} \\
\\ 
\texttt{Έχετε συσσωρεύσει συλλογικά 3 από 8 ποινές.} \\
\texttt{Παρακαλώ δοκιμάστε ξανά!} \\
            }
        }
    }
    & & \\ \\

    \theutterance \stepcounter{utterance}  
    & & \multicolumn{4}{p{0.6\linewidth}}{
        \cellcolor[rgb]{0.9,0.9,0.9}{
            \makecell[{{p{\linewidth}}}]{
                \texttt{\tiny{[P1$\langle$GM]}}
                \texttt{Ποινή: Το (2,2) δεν είναι κενό, αλλά περιέχει '└'.} \\
\\ 
\texttt{Έχετε συσσωρεύσει συλλογικά 3 από 8 ποινές.} \\
\texttt{Παρακαλώ δοκιμάστε ξανά!} \\
            }
        }
    }
    & \\ \\

    \theutterance \stepcounter{utterance}  
    & \multicolumn{4}{p{0.6\linewidth}}{
        \cellcolor[rgb]{0.9,0.9,0.9}{
            \makecell[{{p{\linewidth}}}]{
                \texttt{\tiny{[P1$\rangle$GM]}}
                \texttt{ΠΕΣ: Το (2,2) επίσης κατειλημμένο. Ας πάμε στη σειρά 3: C {-}> (2,3), L {-}> (3,3), P {-}> (4,3). Εντάξει;} \\
            }
        }
    }
    & & \\ \\

    \theutterance \stepcounter{utterance}  
    & & & \multicolumn{2}{p{0.3\linewidth}}{
        \cellcolor[rgb]{0.9,0.9,0.9}{
            \makecell[{{p{\linewidth}}}]{
                \texttt{\tiny{[GM$|$GM]}}
                \texttt{ΠΕΣ: Το (2,2) επίσης κατειλημμένο. Ας πάμε στη σειρά 3: C {-}> (2,3), L {-}> (3,3), P {-}> (4,3). Εντάξει;} \\
            }
        }
    }
    & & \\ \\

    \theutterance \stepcounter{utterance}  
    & & \multicolumn{4}{p{0.6\linewidth}}{
        \cellcolor[rgb]{0.9,0.9,0.9}{
            \makecell[{{p{\linewidth}}}]{
                \texttt{\tiny{[GM$\rangle$P2]}}
                \texttt{Το μήνυμά σας μεταφέρθηκε στον άλλο παίκτη.} \\
\\ 
\\ 
\texttt{Αυτή τη στιγμή παίζετε τον γύρο 4 από συνολικά 12.} \\
\texttt{Έχετε συσσωρεύσει συλλογικά 3 από 8 ποινές.} \\
\texttt{Ο άλλος παίκτης έστειλε αυτό το μήνυμα:} \\
\\ 
\\ 
\texttt{"Το (2,2) επίσης κατειλημμένο. Ας πάμε στη σειρά 3: C {-}> (2,3), L {-}> (3,3), P {-}> (4,3). Εντάξει;"} \\
\\ 
\\ 
\texttt{Ποια είναι η επόμενη εντολή σας;} \\
            }
        }
    }
    & \\ \\

    \theutterance \stepcounter{utterance}  
    & & & \multicolumn{4}{p{0.6\linewidth}}{
        \cellcolor[rgb]{0.9,0.9,0.9}{
            \makecell[{{p{\linewidth}}}]{
                \texttt{\tiny{[GM$\langle$P2]}}
                \texttt{ΠΕΣ: Το (3,3) έχει τοίχιο (από τη θέση (2,2)). Προτείνω τη σειρά 5: C {-}> (2,5), L {-}> (3,5), P {-}> (4,5). Είναι ελεύθερη;} \\
            }
        }
    }
     \\ \\

    \theutterance \stepcounter{utterance}  
    & & & \multicolumn{2}{p{0.3\linewidth}}{
        \cellcolor[rgb]{0.9,0.9,0.9}{
            \makecell[{{p{\linewidth}}}]{
                \texttt{\tiny{[GM$|$GM]}}
                \texttt{ΠΕΣ: Το (3,3) έχει τοίχιο (από τη θέση (2,2)). Προτείνω τη σειρά 5: C {-}> (2,5), L {-}> (3,5), P {-}> (4,5). Είναι ελεύθερη;} \\
            }
        }
    }
    & & \\ \\

    \theutterance \stepcounter{utterance}  
    & & \multicolumn{4}{p{0.6\linewidth}}{
        \cellcolor[rgb]{0.9,0.9,0.9}{
            \makecell[{{p{\linewidth}}}]{
                \texttt{\tiny{[P1$\langle$GM]}}
                \texttt{Το μήνυμά σας μεταφέρθηκε στον άλλο παίκτη.} \\
\\ 
\\ 
\texttt{Αυτή τη στιγμή παίζετε τον γύρο 4 από συνολικά 12.} \\
\texttt{Έχετε συσσωρεύσει συλλογικά 3 από 8 ποινές.} \\
\texttt{Ο άλλος παίκτης έστειλε αυτό το μήνυμα:} \\
\\ 
\\ 
\texttt{"Το (3,3) έχει τοίχιο (από τη θέση (2,2)). Προτείνω τη σειρά 5: C {-}> (2,5), L {-}> (3,5), P {-}> (4,5). Είναι ελεύθερη;"} \\
\\ 
\\ 
\texttt{Ποια είναι η επόμενη εντολή σας;} \\
            }
        }
    }
    & \\ \\

    \theutterance \stepcounter{utterance}  
    & \multicolumn{4}{p{0.6\linewidth}}{
        \cellcolor[rgb]{0.9,0.9,0.9}{
            \makecell[{{p{\linewidth}}}]{
                \texttt{\tiny{[P1$\rangle$GM]}}
                \texttt{ΠΕΣ: Η σειρά 5 στο δικό μου πλέγμα είναι ελεύθερη στα (2,5), (3,5), (4,5). Συμφωνώ: C {-}> (2,5), L {-}> (3,5), P {-}> (4,5). Ξεκινάω εγώ με το C.} \\
            }
        }
    }
    & & \\ \\

    \theutterance \stepcounter{utterance}  
    & & & \multicolumn{2}{p{0.3\linewidth}}{
        \cellcolor[rgb]{0.9,0.9,0.9}{
            \makecell[{{p{\linewidth}}}]{
                \texttt{\tiny{[GM$|$GM]}}
                \texttt{ΠΕΣ: Η σειρά 5 στο δικό μου πλέγμα είναι ελεύθερη στα (2,5), (3,5), (4,5). Συμφωνώ: C {-}> (2,5), L {-}> (3,5), P {-}> (4,5). Ξεκινάω εγώ με το C.} \\
            }
        }
    }
    & & \\ \\

    \theutterance \stepcounter{utterance}  
    & & \multicolumn{4}{p{0.6\linewidth}}{
        \cellcolor[rgb]{0.9,0.9,0.9}{
            \makecell[{{p{\linewidth}}}]{
                \texttt{\tiny{[GM$\rangle$P2]}}
                \texttt{Το μήνυμά σας μεταφέρθηκε στον άλλο παίκτη.} \\
\\ 
\\ 
\texttt{Αυτή τη στιγμή παίζετε τον γύρο 5 από συνολικά 12.} \\
\texttt{Έχετε συσσωρεύσει συλλογικά 3 από 8 ποινές.} \\
\texttt{Ο άλλος παίκτης έστειλε αυτό το μήνυμα:} \\
\\ 
\\ 
\texttt{"Η σειρά 5 στο δικό μου πλέγμα είναι ελεύθερη στα (2,5), (3,5), (4,5). Συμφωνώ: C {-}> (2,5), L {-}> (3,5), P {-}> (4,5). Ξεκινάω εγώ με το C."} \\
\\ 
\\ 
\texttt{Ποια είναι η επόμενη εντολή σας;} \\
            }
        }
    }
    & \\ \\

    \theutterance \stepcounter{utterance}  
    & & & \multicolumn{4}{p{0.6\linewidth}}{
        \cellcolor[rgb]{0.9,0.9,0.9}{
            \makecell[{{p{\linewidth}}}]{
                \texttt{\tiny{[GM$\langle$P2]}}
                \texttt{ΚΙΝΗΣΗ: C, (2,5)} \\
            }
        }
    }
     \\ \\

    \theutterance \stepcounter{utterance}  
    & & & \multicolumn{2}{p{0.3\linewidth}}{
        \cellcolor[rgb]{0.9,0.9,0.9}{
            \makecell[{{p{\linewidth}}}]{
                \texttt{\tiny{[GM$|$GM]}}
                \texttt{ΚΙΝΗΣΗ: C, (2,5)} \\
            }
        }
    }
    & & \\ \\

    \theutterance \stepcounter{utterance}  
    & & & \multicolumn{2}{p{0.3\linewidth}}{
        \cellcolor[rgb]{0.9,0.9,0.9}{
            \makecell[{{p{\linewidth}}}]{
                \texttt{\tiny{[GM$|$GM]}}
                \texttt{Ποινή: Το (2,5) δεν είναι κενό, αλλά περιέχει '│'.} \\
\\ 
\texttt{Έχετε συσσωρεύσει συλλογικά 4 από 8 ποινές.} \\
\texttt{Παρακαλώ δοκιμάστε ξανά!} \\
            }
        }
    }
    & & \\ \\

    \theutterance \stepcounter{utterance}  
    & & \multicolumn{4}{p{0.6\linewidth}}{
        \cellcolor[rgb]{0.9,0.9,0.9}{
            \makecell[{{p{\linewidth}}}]{
                \texttt{\tiny{[GM$\rangle$P2]}}
                \texttt{Ποινή: Το (2,5) δεν είναι κενό, αλλά περιέχει '│'.} \\
\\ 
\texttt{Έχετε συσσωρεύσει συλλογικά 4 από 8 ποινές.} \\
\texttt{Παρακαλώ δοκιμάστε ξανά!} \\
            }
        }
    }
    & \\ \\

    \theutterance \stepcounter{utterance}  
    & & & \multicolumn{4}{p{0.6\linewidth}}{
        \cellcolor[rgb]{0.9,0.9,0.9}{
            \makecell[{{p{\linewidth}}}]{
                \texttt{\tiny{[GM$\langle$P2]}}
                \texttt{ΠΕΣ: Το (2,5) έχει τοίχιο στην πλευρά μου. Ας πάμε στη σειρά 6: C {-}> (2,6), L {-}> (3,6), P {-}> (4,6). Εντάξει;} \\
            }
        }
    }
     \\ \\

    \theutterance \stepcounter{utterance}  
    & & & \multicolumn{2}{p{0.3\linewidth}}{
        \cellcolor[rgb]{0.9,0.9,0.9}{
            \makecell[{{p{\linewidth}}}]{
                \texttt{\tiny{[GM$|$GM]}}
                \texttt{ΠΕΣ: Το (2,5) έχει τοίχιο στην πλευρά μου. Ας πάμε στη σειρά 6: C {-}> (2,6), L {-}> (3,6), P {-}> (4,6). Εντάξει;} \\
            }
        }
    }
    & & \\ \\

    \theutterance \stepcounter{utterance}  
    & & \multicolumn{4}{p{0.6\linewidth}}{
        \cellcolor[rgb]{0.9,0.9,0.9}{
            \makecell[{{p{\linewidth}}}]{
                \texttt{\tiny{[P1$\langle$GM]}}
                \texttt{Το μήνυμά σας μεταφέρθηκε στον άλλο παίκτη.} \\
\\ 
\\ 
\texttt{Αυτή τη στιγμή παίζετε τον γύρο 5 από συνολικά 12.} \\
\texttt{Έχετε συσσωρεύσει συλλογικά 4 από 8 ποινές.} \\
\texttt{Ο άλλος παίκτης έστειλε αυτό το μήνυμα:} \\
\\ 
\\ 
\texttt{"Το (2,5) έχει τοίχιο στην πλευρά μου. Ας πάμε στη σειρά 6: C {-}> (2,6), L {-}> (3,6), P {-}> (4,6). Εντάξει;"} \\
\\ 
\\ 
\texttt{Ποια είναι η επόμενη εντολή σας;} \\
            }
        }
    }
    & \\ \\

    \theutterance \stepcounter{utterance}  
    & \multicolumn{4}{p{0.6\linewidth}}{
        \cellcolor[rgb]{0.9,0.9,0.9}{
            \makecell[{{p{\linewidth}}}]{
                \texttt{\tiny{[P1$\rangle$GM]}}
                \texttt{ΠΕΣ: Η σειρά 6 είναι ελεύθερη εδώ. Συμφωνώ: C {-}> (2,6), L {-}> (3,6), P {-}> (4,6). Ξεκινάω με το C.} \\
            }
        }
    }
    & & \\ \\

    \theutterance \stepcounter{utterance}  
    & & & \multicolumn{2}{p{0.3\linewidth}}{
        \cellcolor[rgb]{0.9,0.9,0.9}{
            \makecell[{{p{\linewidth}}}]{
                \texttt{\tiny{[GM$|$GM]}}
                \texttt{ΠΕΣ: Η σειρά 6 είναι ελεύθερη εδώ. Συμφωνώ: C {-}> (2,6), L {-}> (3,6), P {-}> (4,6). Ξεκινάω με το C.} \\
            }
        }
    }
    & & \\ \\

    \theutterance \stepcounter{utterance}  
    & & \multicolumn{4}{p{0.6\linewidth}}{
        \cellcolor[rgb]{0.9,0.9,0.9}{
            \makecell[{{p{\linewidth}}}]{
                \texttt{\tiny{[GM$\rangle$P2]}}
                \texttt{Το μήνυμά σας μεταφέρθηκε στον άλλο παίκτη.} \\
\\ 
\\ 
\texttt{Αυτή τη στιγμή παίζετε τον γύρο 6 από συνολικά 12.} \\
\texttt{Έχετε συσσωρεύσει συλλογικά 4 από 8 ποινές.} \\
\texttt{Ο άλλος παίκτης έστειλε αυτό το μήνυμα:} \\
\\ 
\\ 
\texttt{"Η σειρά 6 είναι ελεύθερη εδώ. Συμφωνώ: C {-}> (2,6), L {-}> (3,6), P {-}> (4,6). Ξεκινάω με το C."} \\
\\ 
\\ 
\texttt{Ποια είναι η επόμενη εντολή σας;} \\
            }
        }
    }
    & \\ \\

    \theutterance \stepcounter{utterance}  
    & & & \multicolumn{4}{p{0.6\linewidth}}{
        \cellcolor[rgb]{0.9,0.9,0.9}{
            \makecell[{{p{\linewidth}}}]{
                \texttt{\tiny{[GM$\langle$P2]}}
                \texttt{ΚΙΝΗΣΗ: C, (2,6)} \\
            }
        }
    }
     \\ \\

    \theutterance \stepcounter{utterance}  
    & & & \multicolumn{2}{p{0.3\linewidth}}{
        \cellcolor[rgb]{0.9,0.9,0.9}{
            \makecell[{{p{\linewidth}}}]{
                \texttt{\tiny{[GM$|$GM]}}
                \texttt{ΚΙΝΗΣΗ: C, (2,6)} \\
            }
        }
    }
    & & \\ \\

    \theutterance \stepcounter{utterance}  
    & & & \multicolumn{2}{p{0.3\linewidth}}{
        \cellcolor[rgb]{0.9,0.9,0.9}{
            \makecell[{{p{\linewidth}}}]{
                \texttt{\tiny{[GM$|$GM]}}
                \texttt{Ποινή: Το (2,6) δεν είναι κενό, αλλά περιέχει '┼'.} \\
\\ 
\texttt{Έχετε συσσωρεύσει συλλογικά 5 από 8 ποινές.} \\
\texttt{Παρακαλώ δοκιμάστε ξανά!} \\
            }
        }
    }
    & & \\ \\

    \theutterance \stepcounter{utterance}  
    & & \multicolumn{4}{p{0.6\linewidth}}{
        \cellcolor[rgb]{0.9,0.9,0.9}{
            \makecell[{{p{\linewidth}}}]{
                \texttt{\tiny{[GM$\rangle$P2]}}
                \texttt{Ποινή: Το (2,6) δεν είναι κενό, αλλά περιέχει '┼'.} \\
\\ 
\texttt{Έχετε συσσωρεύσει συλλογικά 5 από 8 ποινές.} \\
\texttt{Παρακαλώ δοκιμάστε ξανά!} \\
            }
        }
    }
    & \\ \\

    \theutterance \stepcounter{utterance}  
    & & & \multicolumn{4}{p{0.6\linewidth}}{
        \cellcolor[rgb]{0.9,0.9,0.9}{
            \makecell[{{p{\linewidth}}}]{
                \texttt{\tiny{[GM$\langle$P2]}}
                \texttt{ΠΕΣ: Το (2,6) επίσης κατειλημμένο εδώ. Προτείνω στήλη 5, σειρά 4: C {-}> (5,4), L {-}> (6,4), P {-}> (7,4). Ελεύθερα εδώ. Εντάξει;} \\
            }
        }
    }
     \\ \\

    \theutterance \stepcounter{utterance}  
    & & & \multicolumn{2}{p{0.3\linewidth}}{
        \cellcolor[rgb]{0.9,0.9,0.9}{
            \makecell[{{p{\linewidth}}}]{
                \texttt{\tiny{[GM$|$GM]}}
                \texttt{ΠΕΣ: Το (2,6) επίσης κατειλημμένο εδώ. Προτείνω στήλη 5, σειρά 4: C {-}> (5,4), L {-}> (6,4), P {-}> (7,4). Ελεύθερα εδώ. Εντάξει;} \\
            }
        }
    }
    & & \\ \\

    \theutterance \stepcounter{utterance}  
    & & \multicolumn{4}{p{0.6\linewidth}}{
        \cellcolor[rgb]{0.9,0.9,0.9}{
            \makecell[{{p{\linewidth}}}]{
                \texttt{\tiny{[P1$\langle$GM]}}
                \texttt{Το μήνυμά σας μεταφέρθηκε στον άλλο παίκτη.} \\
\\ 
\\ 
\texttt{Αυτή τη στιγμή παίζετε τον γύρο 6 από συνολικά 12.} \\
\texttt{Έχετε συσσωρεύσει συλλογικά 5 από 8 ποινές.} \\
\texttt{Ο άλλος παίκτης έστειλε αυτό το μήνυμα:} \\
\\ 
\\ 
\texttt{"Το (2,6) επίσης κατειλημμένο εδώ. Προτείνω στήλη 5, σειρά 4: C {-}> (5,4), L {-}> (6,4), P {-}> (7,4). Ελεύθερα εδώ. Εντάξει;"} \\
\\ 
\\ 
\texttt{Ποια είναι η επόμενη εντολή σας;} \\
            }
        }
    }
    & \\ \\

    \theutterance \stepcounter{utterance}  
    & \multicolumn{4}{p{0.6\linewidth}}{
        \cellcolor[rgb]{0.9,0.9,0.9}{
            \makecell[{{p{\linewidth}}}]{
                \texttt{\tiny{[P1$\rangle$GM]}}
                \texttt{ΠΕΣ: Το (5,4) εδώ έχει τοίχιο ('┐'). Δεν λειτουργεί. } \\
\\ 
\texttt{Ας δοκιμάσουμε τη στήλη 5, σειρά 6: C {-}> (5,6), L {-}> (6,6), P {-}> (7,6). Ελεύθερα εδώ. Εντάξει;} \\
            }
        }
    }
    & & \\ \\

    \theutterance \stepcounter{utterance}  
    & & & \multicolumn{2}{p{0.3\linewidth}}{
        \cellcolor[rgb]{0.9,0.9,0.9}{
            \makecell[{{p{\linewidth}}}]{
                \texttt{\tiny{[GM$|$GM]}}
                \texttt{ΠΕΣ: Το (5,4) εδώ έχει τοίχιο ('┐'). Δεν λειτουργεί. } \\
\\ 
\texttt{Ας δοκιμάσουμε τη στήλη 5, σειρά 6: C {-}> (5,6), L {-}> (6,6), P {-}> (7,6). Ελεύθερα εδώ. Εντάξει;} \\
            }
        }
    }
    & & \\ \\

    \theutterance \stepcounter{utterance}  
    & & & \multicolumn{2}{p{0.3\linewidth}}{
        \cellcolor[rgb]{0.9,0.9,0.9}{
            \makecell[{{p{\linewidth}}}]{
                \texttt{\tiny{[GM$|$GM]}}
                \texttt{Invalid format: tail is not empty: '} \\
\\ 
\texttt{Ας δοκιμάσουμε τη στήλη 5, σειρά 6: C {-}> (5,6), L {-}> (6,6), P {-}> (7,6). Ελεύθερα εδώ. Εντάξει;} \\
            }
        }
    }
    & & \\ \\

    \theutterance \stepcounter{utterance}  
    & & \multicolumn{4}{p{0.6\linewidth}}{
        \cellcolor[rgb]{0.9,0.9,0.9}{
            \makecell[{{p{\linewidth}}}]{
                \texttt{\tiny{[P1$\langle$GM]}}
                \texttt{Ποινή: Το μήνυμά σας δεν πρέπει να περιέχει τίποτα μετά την εντολή!} \\
\\ 
\texttt{Βεβαιωθείτε ότι η απάντησή σας περιέχει μόνο `ΠΕΣ: <ΜΉΝΥΜΑ>` ή `ΚΙΝΗΣΗ: <ΑΝΤΙΚΕΊΜΕΝΟ>, (<X>, <Y>)`, και τίποτα άλλο!} \\
\\ 
\texttt{Έχετε συσσωρεύσει συλλογικά 6 από 8 ποινές.} \\
\texttt{Παρακαλώ δοκιμάστε ξανά!} \\
            }
        }
    }
    & \\ \\

    \theutterance \stepcounter{utterance}  
    & \multicolumn{4}{p{0.6\linewidth}}{
        \cellcolor[rgb]{0.9,0.9,0.9}{
            \makecell[{{p{\linewidth}}}]{
                \texttt{\tiny{[P1$\rangle$GM]}}
                \texttt{ΠΕΣ: Το (5,4) εδώ έχει τοίχιο. Δεν λειτουργεί. Ας δοκιμάσουμε στήλη 5, σειρά 6: C {-}> (5,6), L {-}> (6,6), P {-}> (7,6). Ελεύθερα εδώ. Εντάξει;} \\
            }
        }
    }
    & & \\ \\

    \theutterance \stepcounter{utterance}  
    & & & \multicolumn{2}{p{0.3\linewidth}}{
        \cellcolor[rgb]{0.9,0.9,0.9}{
            \makecell[{{p{\linewidth}}}]{
                \texttt{\tiny{[GM$|$GM]}}
                \texttt{ΠΕΣ: Το (5,4) εδώ έχει τοίχιο. Δεν λειτουργεί. Ας δοκιμάσουμε στήλη 5, σειρά 6: C {-}> (5,6), L {-}> (6,6), P {-}> (7,6). Ελεύθερα εδώ. Εντάξει;} \\
            }
        }
    }
    & & \\ \\

    \theutterance \stepcounter{utterance}  
    & & \multicolumn{4}{p{0.6\linewidth}}{
        \cellcolor[rgb]{0.9,0.9,0.9}{
            \makecell[{{p{\linewidth}}}]{
                \texttt{\tiny{[GM$\rangle$P2]}}
                \texttt{Το μήνυμά σας μεταφέρθηκε στον άλλο παίκτη.} \\
\\ 
\\ 
\texttt{Αυτή τη στιγμή παίζετε τον γύρο 7 από συνολικά 12.} \\
\texttt{Έχετε συσσωρεύσει συλλογικά 6 από 8 ποινές.} \\
\texttt{Ο άλλος παίκτης έστειλε αυτό το μήνυμα:} \\
\\ 
\\ 
\texttt{"Το (5,4) εδώ έχει τοίχιο. Δεν λειτουργεί. Ας δοκιμάσουμε στήλη 5, σειρά 6: C {-}> (5,6), L {-}> (6,6), P {-}> (7,6). Ελεύθερα εδώ. Εντάξει;"} \\
\\ 
\\ 
\texttt{Ποια είναι η επόμενη εντολή σας;} \\
            }
        }
    }
    & \\ \\

    \theutterance \stepcounter{utterance}  
    & & & \multicolumn{4}{p{0.6\linewidth}}{
        \cellcolor[rgb]{0.9,0.9,0.9}{
            \makecell[{{p{\linewidth}}}]{
                \texttt{\tiny{[GM$\langle$P2]}}
                \texttt{ΠΕΣ: Έχεις δίκιο, το (5,6) εδώ είναι τοίχιο. Το (6,6) και (7,6) είναι ελεύθερα. Το (5,4) εδώ είναι ελεύθερο. Ας βάλουμε: C {-}> (5,4), L {-}> (6,4), P {-}> (7,4). Σου φαίνεται καλό;} \\
            }
        }
    }
     \\ \\

    \theutterance \stepcounter{utterance}  
    & & & \multicolumn{2}{p{0.3\linewidth}}{
        \cellcolor[rgb]{0.9,0.9,0.9}{
            \makecell[{{p{\linewidth}}}]{
                \texttt{\tiny{[GM$|$GM]}}
                \texttt{ΠΕΣ: Έχεις δίκιο, το (5,6) εδώ είναι τοίχιο. Το (6,6) και (7,6) είναι ελεύθερα. Το (5,4) εδώ είναι ελεύθερο. Ας βάλουμε: C {-}> (5,4), L {-}> (6,4), P {-}> (7,4). Σου φαίνεται καλό;} \\
            }
        }
    }
    & & \\ \\

    \theutterance \stepcounter{utterance}  
    & & \multicolumn{4}{p{0.6\linewidth}}{
        \cellcolor[rgb]{0.9,0.9,0.9}{
            \makecell[{{p{\linewidth}}}]{
                \texttt{\tiny{[P1$\langle$GM]}}
                \texttt{Το μήνυμά σας μεταφέρθηκε στον άλλο παίκτη.} \\
\\ 
\\ 
\texttt{Αυτή τη στιγμή παίζετε τον γύρο 7 από συνολικά 12.} \\
\texttt{Έχετε συσσωρεύσει συλλογικά 6 από 8 ποινές.} \\
\texttt{Ο άλλος παίκτης έστειλε αυτό το μήνυμα:} \\
\\ 
\\ 
\texttt{"Έχεις δίκιο, το (5,6) εδώ είναι τοίχιο. Το (6,6) και (7,6) είναι ελεύθερα. Το (5,4) εδώ είναι ελεύθερο. Ας βάλουμε: C {-}> (5,4), L {-}> (6,4), P {-}> (7,4). Σου φαίνεται καλό;"} \\
\\ 
\\ 
\texttt{Ποια είναι η επόμενη εντολή σας;} \\
            }
        }
    }
    & \\ \\

    \theutterance \stepcounter{utterance}  
    & \multicolumn{4}{p{0.6\linewidth}}{
        \cellcolor[rgb]{0.9,0.9,0.9}{
            \makecell[{{p{\linewidth}}}]{
                \texttt{\tiny{[P1$\rangle$GM]}}
                \texttt{ΠΕΣ: Το (5,4) εδώ είναι τοίχιο ('┐'), οπότε δεν μπορώ να βάλω τίποτε εκεί. Προτείνω στήλη 5, σειρά 6: C {-}> (5,6), L {-}> (6,6), P {-}> (7,6). Εδώ είναι όλα ελεύθερα. Συμφωνείς;} \\
            }
        }
    }
    & & \\ \\

    \theutterance \stepcounter{utterance}  
    & & & \multicolumn{2}{p{0.3\linewidth}}{
        \cellcolor[rgb]{0.9,0.9,0.9}{
            \makecell[{{p{\linewidth}}}]{
                \texttt{\tiny{[GM$|$GM]}}
                \texttt{ΠΕΣ: Το (5,4) εδώ είναι τοίχιο ('┐'), οπότε δεν μπορώ να βάλω τίποτε εκεί. Προτείνω στήλη 5, σειρά 6: C {-}> (5,6), L {-}> (6,6), P {-}> (7,6). Εδώ είναι όλα ελεύθερα. Συμφωνείς;} \\
            }
        }
    }
    & & \\ \\

    \theutterance \stepcounter{utterance}  
    & & \multicolumn{4}{p{0.6\linewidth}}{
        \cellcolor[rgb]{0.9,0.9,0.9}{
            \makecell[{{p{\linewidth}}}]{
                \texttt{\tiny{[GM$\rangle$P2]}}
                \texttt{Το μήνυμά σας μεταφέρθηκε στον άλλο παίκτη.} \\
\\ 
\\ 
\texttt{Αυτή τη στιγμή παίζετε τον γύρο 8 από συνολικά 12.} \\
\texttt{Έχετε συσσωρεύσει συλλογικά 6 από 8 ποινές.} \\
\texttt{Ο άλλος παίκτης έστειλε αυτό το μήνυμα:} \\
\\ 
\\ 
\texttt{"Το (5,4) εδώ είναι τοίχιο ('┐'), οπότε δεν μπορώ να βάλω τίποτε εκεί. Προτείνω στήλη 5, σειρά 6: C {-}> (5,6), L {-}> (6,6), P {-}> (7,6). Εδώ είναι όλα ελεύθερα. Συμφωνείς;"} \\
\\ 
\\ 
\texttt{Ποια είναι η επόμενη εντολή σας;} \\
            }
        }
    }
    & \\ \\

    \theutterance \stepcounter{utterance}  
    & & & \multicolumn{4}{p{0.6\linewidth}}{
        \cellcolor[rgb]{0.9,0.9,0.9}{
            \makecell[{{p{\linewidth}}}]{
                \texttt{\tiny{[GM$\langle$P2]}}
                \texttt{ΠΕΣ: Συμφωνώ για τη στήλη 5, σειρά 6: C {-}> (5,6), L {-}> (6,6), P {-}> (7,6). Ξεκινάω εγώ με το C.} \\
            }
        }
    }
     \\ \\

    \theutterance \stepcounter{utterance}  
    & & & \multicolumn{2}{p{0.3\linewidth}}{
        \cellcolor[rgb]{0.9,0.9,0.9}{
            \makecell[{{p{\linewidth}}}]{
                \texttt{\tiny{[GM$|$GM]}}
                \texttt{ΠΕΣ: Συμφωνώ για τη στήλη 5, σειρά 6: C {-}> (5,6), L {-}> (6,6), P {-}> (7,6). Ξεκινάω εγώ με το C.} \\
            }
        }
    }
    & & \\ \\

    \theutterance \stepcounter{utterance}  
    & & \multicolumn{4}{p{0.6\linewidth}}{
        \cellcolor[rgb]{0.9,0.9,0.9}{
            \makecell[{{p{\linewidth}}}]{
                \texttt{\tiny{[P1$\langle$GM]}}
                \texttt{Το μήνυμά σας μεταφέρθηκε στον άλλο παίκτη.} \\
\\ 
\\ 
\texttt{Αυτή τη στιγμή παίζετε τον γύρο 8 από συνολικά 12.} \\
\texttt{Έχετε συσσωρεύσει συλλογικά 6 από 8 ποινές.} \\
\texttt{Ο άλλος παίκτης έστειλε αυτό το μήνυμα:} \\
\\ 
\\ 
\texttt{"Συμφωνώ για τη στήλη 5, σειρά 6: C {-}> (5,6), L {-}> (6,6), P {-}> (7,6). Ξεκινάω εγώ με το C."} \\
\\ 
\\ 
\texttt{Ποια είναι η επόμενη εντολή σας;} \\
            }
        }
    }
    & \\ \\

    \theutterance \stepcounter{utterance}  
    & \multicolumn{4}{p{0.6\linewidth}}{
        \cellcolor[rgb]{0.9,0.9,0.9}{
            \makecell[{{p{\linewidth}}}]{
                \texttt{\tiny{[P1$\rangle$GM]}}
                \texttt{ΚΙΝΗΣΗ: C, (5, 6)} \\
            }
        }
    }
    & & \\ \\

    \theutterance \stepcounter{utterance}  
    & & & \multicolumn{2}{p{0.3\linewidth}}{
        \cellcolor[rgb]{0.9,0.9,0.9}{
            \makecell[{{p{\linewidth}}}]{
                \texttt{\tiny{[GM$|$GM]}}
                \texttt{ΚΙΝΗΣΗ: C, (5, 6)} \\
            }
        }
    }
    & & \\ \\

    \theutterance \stepcounter{utterance}  
    & & & \multicolumn{2}{p{0.3\linewidth}}{
        \cellcolor[rgb]{0.9,0.9,0.9}{
            \makecell[{{p{\linewidth}}}]{
                \texttt{\tiny{[GM$|$GM]}}
                \texttt{Ποινή: Το (5,6) δεν είναι κενό, αλλά περιέχει '─'.} \\
\\ 
\texttt{Έχετε συσσωρεύσει συλλογικά 7 από 8 ποινές.} \\
\texttt{Παρακαλώ δοκιμάστε ξανά!} \\
            }
        }
    }
    & & \\ \\

    \theutterance \stepcounter{utterance}  
    & & \multicolumn{4}{p{0.6\linewidth}}{
        \cellcolor[rgb]{0.9,0.9,0.9}{
            \makecell[{{p{\linewidth}}}]{
                \texttt{\tiny{[P1$\langle$GM]}}
                \texttt{Ποινή: Το (5,6) δεν είναι κενό, αλλά περιέχει '─'.} \\
\\ 
\texttt{Έχετε συσσωρεύσει συλλογικά 7 από 8 ποινές.} \\
\texttt{Παρακαλώ δοκιμάστε ξανά!} \\
            }
        }
    }
    & \\ \\

    \theutterance \stepcounter{utterance}  
    & \multicolumn{4}{p{0.6\linewidth}}{
        \cellcolor[rgb]{0.9,0.9,0.9}{
            \makecell[{{p{\linewidth}}}]{
                \texttt{\tiny{[P1$\rangle$GM]}}
                \texttt{ΠΕΣ: Το (5,6) εδώ έχει τοίχιο ('─'). Μια στιγμή, ελέγχω το πλέγμα μου ξανά... } \\
\\ 
\texttt{Οι ελεύθερες θέσεις στη σειρά 6 είναι μόνο το (2,6) και το (7,6). Στη στήλη 5, μόνο το (5,3) είναι ελεύθερο. } \\
\\ 
\texttt{Προτείνω να χρησιμοποιήσουμε τη στήλη 3, σειρά 5: C {-}> (3,5), L {-}> (4,5), P {-}> (5,5). Ελεύθερα εδώ. Εντάξει;} \\
            }
        }
    }
    & & \\ \\

    \theutterance \stepcounter{utterance}  
    & & & \multicolumn{2}{p{0.3\linewidth}}{
        \cellcolor[rgb]{0.9,0.9,0.9}{
            \makecell[{{p{\linewidth}}}]{
                \texttt{\tiny{[GM$|$GM]}}
                \texttt{ΠΕΣ: Το (5,6) εδώ έχει τοίχιο ('─'). Μια στιγμή, ελέγχω το πλέγμα μου ξανά... } \\
\\ 
\texttt{Οι ελεύθερες θέσεις στη σειρά 6 είναι μόνο το (2,6) και το (7,6). Στη στήλη 5, μόνο το (5,3) είναι ελεύθερο. } \\
\\ 
\texttt{Προτείνω να χρησιμοποιήσουμε τη στήλη 3, σειρά 5: C {-}> (3,5), L {-}> (4,5), P {-}> (5,5). Ελεύθερα εδώ. Εντάξει;} \\
            }
        }
    }
    & & \\ \\

    \theutterance \stepcounter{utterance}  
    & & & \multicolumn{2}{p{0.3\linewidth}}{
        \cellcolor[rgb]{0.9,0.9,0.9}{
            \makecell[{{p{\linewidth}}}]{
                \texttt{\tiny{[GM$|$GM]}}
                \texttt{Invalid format: tail is not empty: '} \\
\\ 
\texttt{Οι ελεύθερες θέσεις στη σειρά 6 είναι μόνο το (2,6) και το (7,6). Στη στήλη 5, μόνο το (5,3) είναι ελεύθερο. } \\
\\ 
\texttt{Προτείνω να χρησιμοποιήσουμε τη στήλη 3, σειρά 5: C {-}> (3,5), L {-}> (4,5), P {-}> (5,5). Ελεύθερα εδώ. Εντάξει;} \\
            }
        }
    }
    & & \\ \\

    \theutterance \stepcounter{utterance}  
    & & \multicolumn{4}{p{0.6\linewidth}}{
        \cellcolor[rgb]{0.9,0.9,0.9}{
            \makecell[{{p{\linewidth}}}]{
                \texttt{\tiny{[P1$\langle$GM]}}
                \texttt{Ποινή: Το μήνυμά σας δεν πρέπει να περιέχει τίποτα μετά την εντολή!} \\
\\ 
\texttt{Βεβαιωθείτε ότι η απάντησή σας περιέχει μόνο `ΠΕΣ: <ΜΉΝΥΜΑ>` ή `ΚΙΝΗΣΗ: <ΑΝΤΙΚΕΊΜΕΝΟ>, (<X>, <Y>)`, και τίποτα άλλο!} \\
\\ 
\texttt{Έχετε συσσωρεύσει συλλογικά 8 από 8 ποινές.} \\
\texttt{Παρακαλώ δοκιμάστε ξανά!} \\
            }
        }
    }
    & \\ \\

    \theutterance \stepcounter{utterance}  
    & \multicolumn{4}{p{0.6\linewidth}}{
        \cellcolor[rgb]{0.9,0.9,0.9}{
            \makecell[{{p{\linewidth}}}]{
                \texttt{\tiny{[P1$\rangle$GM]}}
                \texttt{ΠΕΣ: Το (5,6) έχει τοίχιο. Εδώ ελεύθερα: (3,5), (4,5), (5,5). Προτείνω: C {-}> (3,5), L {-}> (4,5), P {-}> (5,5). Εντάξει;} \\
            }
        }
    }
    & & \\ \\

    \theutterance \stepcounter{utterance}  
    & & & \multicolumn{2}{p{0.3\linewidth}}{
        \cellcolor[rgb]{0.9,0.9,0.9}{
            \makecell[{{p{\linewidth}}}]{
                \texttt{\tiny{[GM$|$GM]}}
                \texttt{ΠΕΣ: Το (5,6) έχει τοίχιο. Εδώ ελεύθερα: (3,5), (4,5), (5,5). Προτείνω: C {-}> (3,5), L {-}> (4,5), P {-}> (5,5). Εντάξει;} \\
            }
        }
    }
    & & \\ \\

    \theutterance \stepcounter{utterance}  
    & & \multicolumn{4}{p{0.6\linewidth}}{
        \cellcolor[rgb]{0.9,0.9,0.9}{
            \makecell[{{p{\linewidth}}}]{
                \texttt{\tiny{[GM$\rangle$P2]}}
                \texttt{Το μήνυμά σας μεταφέρθηκε στον άλλο παίκτη.} \\
\\ 
\\ 
\texttt{Αυτή τη στιγμή παίζετε τον γύρο 9 από συνολικά 12.} \\
\texttt{Έχετε συσσωρεύσει συλλογικά 8 από 8 ποινές.} \\
\texttt{Ο άλλος παίκτης έστειλε αυτό το μήνυμα:} \\
\\ 
\\ 
\texttt{"Το (5,6) έχει τοίχιο. Εδώ ελεύθερα: (3,5), (4,5), (5,5). Προτείνω: C {-}> (3,5), L {-}> (4,5), P {-}> (5,5). Εντάξει;"} \\
\\ 
\\ 
\texttt{Ποια είναι η επόμενη εντολή σας;} \\
            }
        }
    }
    & \\ \\

    \theutterance \stepcounter{utterance}  
    & & & \multicolumn{4}{p{0.6\linewidth}}{
        \cellcolor[rgb]{0.9,0.9,0.9}{
            \makecell[{{p{\linewidth}}}]{
                \texttt{\tiny{[GM$\langle$P2]}}
                \texttt{ΠΕΣ: Τα (3,5), (4,5), (5,5) είναι ελεύθερα εδώ. Συμφωνώ: C {-}> (3,5), L {-}> (4,5), P {-}> (5,5). Ξεκινάω με το C.} \\
            }
        }
    }
     \\ \\

    \theutterance \stepcounter{utterance}  
    & & & \multicolumn{2}{p{0.3\linewidth}}{
        \cellcolor[rgb]{0.9,0.9,0.9}{
            \makecell[{{p{\linewidth}}}]{
                \texttt{\tiny{[GM$|$GM]}}
                \texttt{ΠΕΣ: Τα (3,5), (4,5), (5,5) είναι ελεύθερα εδώ. Συμφωνώ: C {-}> (3,5), L {-}> (4,5), P {-}> (5,5). Ξεκινάω με το C.} \\
            }
        }
    }
    & & \\ \\

    \theutterance \stepcounter{utterance}  
    & & \multicolumn{4}{p{0.6\linewidth}}{
        \cellcolor[rgb]{0.9,0.9,0.9}{
            \makecell[{{p{\linewidth}}}]{
                \texttt{\tiny{[P1$\langle$GM]}}
                \texttt{Το μήνυμά σας μεταφέρθηκε στον άλλο παίκτη.} \\
\\ 
\\ 
\texttt{Αυτή τη στιγμή παίζετε τον γύρο 9 από συνολικά 12.} \\
\texttt{Έχετε συσσωρεύσει συλλογικά 8 από 8 ποινές.} \\
\texttt{Ο άλλος παίκτης έστειλε αυτό το μήνυμα:} \\
\\ 
\\ 
\texttt{"Τα (3,5), (4,5), (5,5) είναι ελεύθερα εδώ. Συμφωνώ: C {-}> (3,5), L {-}> (4,5), P {-}> (5,5). Ξεκινάω με το C."} \\
\\ 
\\ 
\texttt{Ποια είναι η επόμενη εντολή σας;} \\
            }
        }
    }
    & \\ \\

    \theutterance \stepcounter{utterance}  
    & \multicolumn{4}{p{0.6\linewidth}}{
        \cellcolor[rgb]{0.9,0.9,0.9}{
            \makecell[{{p{\linewidth}}}]{
                \texttt{\tiny{[P1$\rangle$GM]}}
                \texttt{ΚΙΝΗΣΗ: C, (3, 5)} \\
            }
        }
    }
    & & \\ \\

    \theutterance \stepcounter{utterance}  
    & & & \multicolumn{2}{p{0.3\linewidth}}{
        \cellcolor[rgb]{0.9,0.9,0.9}{
            \makecell[{{p{\linewidth}}}]{
                \texttt{\tiny{[GM$|$GM]}}
                \texttt{ΚΙΝΗΣΗ: C, (3, 5)} \\
            }
        }
    }
    & & \\ \\

    \theutterance \stepcounter{utterance}  
    & & & \multicolumn{2}{p{0.3\linewidth}}{
        \cellcolor[rgb]{0.9,0.9,0.9}{
            \makecell[{{p{\linewidth}}}]{
                \texttt{\tiny{[GM$|$GM]}}
                \texttt{C moved to (3, 5)} \\
            }
        }
    }
    & & \\ \\

    \theutterance \stepcounter{utterance}  
    & & \multicolumn{4}{p{0.6\linewidth}}{
        \cellcolor[rgb]{0.9,0.9,0.9}{
            \makecell[{{p{\linewidth}}}]{
                \texttt{\tiny{[GM$\rangle$P2]}}
                \texttt{Το μήνυμά σας μεταφέρθηκε στον άλλο παίκτη.} \\
\\ 
\\ 
\texttt{Αυτή τη στιγμή παίζετε τον γύρο 10 από συνολικά 12.} \\
\texttt{Έχετε συσσωρεύσει συλλογικά 8 από 8 ποινές.} \\
\texttt{Ο άλλος παίκτης μετακίνησε ένα αντικείμενο στο πλέγμα του.} \\
\\ 
\\ 
\texttt{Ποια είναι η επόμενη εντολή σας;} \\
            }
        }
    }
    & \\ \\

    \theutterance \stepcounter{utterance}  
    & & & \multicolumn{4}{p{0.6\linewidth}}{
        \cellcolor[rgb]{0.9,0.9,0.9}{
            \makecell[{{p{\linewidth}}}]{
                \texttt{\tiny{[GM$\langle$P2]}}
                \texttt{ΚΙΝΗΣΗ: C, (3,5)} \\
            }
        }
    }
     \\ \\

    \theutterance \stepcounter{utterance}  
    & & & \multicolumn{2}{p{0.3\linewidth}}{
        \cellcolor[rgb]{0.9,0.9,0.9}{
            \makecell[{{p{\linewidth}}}]{
                \texttt{\tiny{[GM$|$GM]}}
                \texttt{ΚΙΝΗΣΗ: C, (3,5)} \\
            }
        }
    }
    & & \\ \\

    \theutterance \stepcounter{utterance}  
    & & & \multicolumn{2}{p{0.3\linewidth}}{
        \cellcolor[rgb]{0.9,0.9,0.9}{
            \makecell[{{p{\linewidth}}}]{
                \texttt{\tiny{[GM$|$GM]}}
                \texttt{C moved to (3, 5)} \\
            }
        }
    }
    & & \\ \\

    \theutterance \stepcounter{utterance}  
    & & \multicolumn{4}{p{0.6\linewidth}}{
        \cellcolor[rgb]{0.9,0.9,0.9}{
            \makecell[{{p{\linewidth}}}]{
                \texttt{\tiny{[P1$\langle$GM]}}
                \texttt{Μεταφέρθηκε επιτυχώς το 'C' στο (3,5). Το ενημερωμένο πλέγμα σας μοιάζει ως εξής:} \\
\texttt{```} \\
\texttt{ 1234567} \\
\texttt{╔═╤══╤══╗} \\
\texttt{║◌│◌◌│◌◌║ 1} \\
\texttt{║◌└┬─┴──╢ 2} \\
\texttt{║◌◌│L◌◌◌║ 3} \\
\texttt{╟─┬┘◌◌┌─╢ 4} \\
\texttt{║◌│C◌◌│◌║ 5} \\
\texttt{╟─┼─┬─┤P║ 6} \\
\texttt{║◌│◌│◌│◌║ 7} \\
\texttt{╚═╧═╧═╧═╝} \\
\texttt{```} \\
\texttt{Αυτή τη στιγμή παίζετε τον γύρο 10 από συνολικά 12.} \\
\texttt{Έχετε συσσωρεύσει συλλογικά 8 από 8 ποινές.} \\
\texttt{Ο άλλος παίκτης μετακίνησε ένα αντικείμενο στο πλέγμα του.} \\
\\ 
\\ 
\texttt{Ποια είναι η επόμενη εντολή σας;} \\
            }
        }
    }
    & \\ \\

    \theutterance \stepcounter{utterance}  
    & \multicolumn{4}{p{0.6\linewidth}}{
        \cellcolor[rgb]{0.9,0.9,0.9}{
            \makecell[{{p{\linewidth}}}]{
                \texttt{\tiny{[P1$\rangle$GM]}}
                \texttt{ΚΙΝΗΣΗ: L, (4, 5)} \\
            }
        }
    }
    & & \\ \\

    \theutterance \stepcounter{utterance}  
    & & & \multicolumn{2}{p{0.3\linewidth}}{
        \cellcolor[rgb]{0.9,0.9,0.9}{
            \makecell[{{p{\linewidth}}}]{
                \texttt{\tiny{[GM$|$GM]}}
                \texttt{ΚΙΝΗΣΗ: L, (4, 5)} \\
            }
        }
    }
    & & \\ \\

    \theutterance \stepcounter{utterance}  
    & & & \multicolumn{2}{p{0.3\linewidth}}{
        \cellcolor[rgb]{0.9,0.9,0.9}{
            \makecell[{{p{\linewidth}}}]{
                \texttt{\tiny{[GM$|$GM]}}
                \texttt{L moved to (4, 5)} \\
            }
        }
    }
    & & \\ \\

    \theutterance \stepcounter{utterance}  
    & & \multicolumn{4}{p{0.6\linewidth}}{
        \cellcolor[rgb]{0.9,0.9,0.9}{
            \makecell[{{p{\linewidth}}}]{
                \texttt{\tiny{[GM$\rangle$P2]}}
                \texttt{Μεταφέρθηκε επιτυχώς το 'C' στο (3,5). Το ενημερωμένο πλέγμα σας μοιάζει ως εξής:} \\
\texttt{```} \\
\texttt{ 1234567} \\
\texttt{╔═╤══╤══╗} \\
\texttt{║◌│◌◌│◌◌║ 1} \\
\texttt{║L└┬─┴──╢ 2} \\
\texttt{║◌◌│◌◌◌◌║ 3} \\
\texttt{╟─┬┘◌◌┌─╢ 4} \\
\texttt{║◌│C◌◌│◌║ 5} \\
\texttt{╟─┼─┬─┤◌║ 6} \\
\texttt{║◌│P│◌│◌║ 7} \\
\texttt{╚═╧═╧═╧═╝} \\
\texttt{```} \\
\texttt{Αυτή τη στιγμή παίζετε τον γύρο 11 από συνολικά 12.} \\
\texttt{Έχετε συσσωρεύσει συλλογικά 8 από 8 ποινές.} \\
\texttt{Ο άλλος παίκτης μετακίνησε ένα αντικείμενο στο πλέγμα του.} \\
\\ 
\\ 
\texttt{Ποια είναι η επόμενη εντολή σας;} \\
            }
        }
    }
    & \\ \\

    \theutterance \stepcounter{utterance}  
    & & & \multicolumn{4}{p{0.6\linewidth}}{
        \cellcolor[rgb]{0.9,0.9,0.9}{
            \makecell[{{p{\linewidth}}}]{
                \texttt{\tiny{[GM$\langle$P2]}}
                \texttt{ΚΙΝΗΣΗ: L, (4,5)} \\
            }
        }
    }
     \\ \\

    \theutterance \stepcounter{utterance}  
    & & & \multicolumn{2}{p{0.3\linewidth}}{
        \cellcolor[rgb]{0.9,0.9,0.9}{
            \makecell[{{p{\linewidth}}}]{
                \texttt{\tiny{[GM$|$GM]}}
                \texttt{ΚΙΝΗΣΗ: L, (4,5)} \\
            }
        }
    }
    & & \\ \\

    \theutterance \stepcounter{utterance}  
    & & & \multicolumn{2}{p{0.3\linewidth}}{
        \cellcolor[rgb]{0.9,0.9,0.9}{
            \makecell[{{p{\linewidth}}}]{
                \texttt{\tiny{[GM$|$GM]}}
                \texttt{L moved to (4, 5)} \\
            }
        }
    }
    & & \\ \\

    \theutterance \stepcounter{utterance}  
    & & \multicolumn{4}{p{0.6\linewidth}}{
        \cellcolor[rgb]{0.9,0.9,0.9}{
            \makecell[{{p{\linewidth}}}]{
                \texttt{\tiny{[P1$\langle$GM]}}
                \texttt{Μεταφέρθηκε επιτυχώς το 'L' στο (4,5). Το ενημερωμένο πλέγμα σας μοιάζει ως εξής:} \\
\texttt{```} \\
\texttt{ 1234567} \\
\texttt{╔═╤══╤══╗} \\
\texttt{║◌│◌◌│◌◌║ 1} \\
\texttt{║◌└┬─┴──╢ 2} \\
\texttt{║◌◌│◌◌◌◌║ 3} \\
\texttt{╟─┬┘◌◌┌─╢ 4} \\
\texttt{║◌│CL◌│◌║ 5} \\
\texttt{╟─┼─┬─┤P║ 6} \\
\texttt{║◌│◌│◌│◌║ 7} \\
\texttt{╚═╧═╧═╧═╝} \\
\texttt{```} \\
\texttt{Αυτή τη στιγμή παίζετε τον γύρο 11 από συνολικά 12.} \\
\texttt{Έχετε συσσωρεύσει συλλογικά 8 από 8 ποινές.} \\
\texttt{Ο άλλος παίκτης μετακίνησε ένα αντικείμενο στο πλέγμα του.} \\
\\ 
\\ 
\texttt{Ποια είναι η επόμενη εντολή σας;} \\
            }
        }
    }
    & \\ \\

    \theutterance \stepcounter{utterance}  
    & \multicolumn{4}{p{0.6\linewidth}}{
        \cellcolor[rgb]{0.9,0.9,0.9}{
            \makecell[{{p{\linewidth}}}]{
                \texttt{\tiny{[P1$\rangle$GM]}}
                \texttt{ΚΙΝΗΣΗ: P, (5, 5)} \\
            }
        }
    }
    & & \\ \\

    \theutterance \stepcounter{utterance}  
    & & & \multicolumn{2}{p{0.3\linewidth}}{
        \cellcolor[rgb]{0.9,0.9,0.9}{
            \makecell[{{p{\linewidth}}}]{
                \texttt{\tiny{[GM$|$GM]}}
                \texttt{ΚΙΝΗΣΗ: P, (5, 5)} \\
            }
        }
    }
    & & \\ \\

    \theutterance \stepcounter{utterance}  
    & & & \multicolumn{2}{p{0.3\linewidth}}{
        \cellcolor[rgb]{0.9,0.9,0.9}{
            \makecell[{{p{\linewidth}}}]{
                \texttt{\tiny{[GM$|$GM]}}
                \texttt{P moved to (5, 5)} \\
            }
        }
    }
    & & \\ \\

    \theutterance \stepcounter{utterance}  
    & & \multicolumn{4}{p{0.6\linewidth}}{
        \cellcolor[rgb]{0.9,0.9,0.9}{
            \makecell[{{p{\linewidth}}}]{
                \texttt{\tiny{[GM$\rangle$P2]}}
                \texttt{Μεταφέρθηκε επιτυχώς το 'L' στο (4,5). Το ενημερωμένο πλέγμα σας μοιάζει ως εξής:} \\
\texttt{```} \\
\texttt{ 1234567} \\
\texttt{╔═╤══╤══╗} \\
\texttt{║◌│◌◌│◌◌║ 1} \\
\texttt{║◌└┬─┴──╢ 2} \\
\texttt{║◌◌│◌◌◌◌║ 3} \\
\texttt{╟─┬┘◌◌┌─╢ 4} \\
\texttt{║◌│CL◌│◌║ 5} \\
\texttt{╟─┼─┬─┤◌║ 6} \\
\texttt{║◌│P│◌│◌║ 7} \\
\texttt{╚═╧═╧═╧═╝} \\
\texttt{```} \\
\texttt{Αυτή τη στιγμή παίζετε τον γύρο 12 από συνολικά 12.} \\
\texttt{Έχετε συσσωρεύσει συλλογικά 8 από 8 ποινές.} \\
\texttt{Ο άλλος παίκτης μετακίνησε ένα αντικείμενο στο πλέγμα του.} \\
\\ 
\\ 
\texttt{Ποια είναι η επόμενη εντολή σας;} \\
            }
        }
    }
    & \\ \\

    \theutterance \stepcounter{utterance}  
    & & & \multicolumn{4}{p{0.6\linewidth}}{
        \cellcolor[rgb]{0.9,0.9,0.9}{
            \makecell[{{p{\linewidth}}}]{
                \texttt{\tiny{[GM$\langle$P2]}}
                \texttt{ΚΙΝΗΣΗ: P, (5,5)} \\
            }
        }
    }
     \\ \\

    \theutterance \stepcounter{utterance}  
    & & & \multicolumn{2}{p{0.3\linewidth}}{
        \cellcolor[rgb]{0.9,0.9,0.9}{
            \makecell[{{p{\linewidth}}}]{
                \texttt{\tiny{[GM$|$GM]}}
                \texttt{ΚΙΝΗΣΗ: P, (5,5)} \\
            }
        }
    }
    & & \\ \\

    \theutterance \stepcounter{utterance}  
    & & & \multicolumn{2}{p{0.3\linewidth}}{
        \cellcolor[rgb]{0.9,0.9,0.9}{
            \makecell[{{p{\linewidth}}}]{
                \texttt{\tiny{[GM$|$GM]}}
                \texttt{P moved to (5, 5)} \\
            }
        }
    }
    & & \\ \\

    \theutterance \stepcounter{utterance}  
    & & & \multicolumn{2}{p{0.3\linewidth}}{
        \cellcolor[rgb]{0.9,0.9,0.9}{
            \makecell[{{p{\linewidth}}}]{
                \texttt{\tiny{[GM$|$GM]}}
                \texttt{Maximum number of rounds reached} \\
            }
        }
    }
    & & \\ \\

    \theutterance \stepcounter{utterance}  
    & & & \multicolumn{2}{p{0.3\linewidth}}{
        \cellcolor[rgb]{0.9,0.9,0.9}{
            \makecell[{{p{\linewidth}}}]{
                \texttt{\tiny{[GM$|$GM]}}
                \texttt{Initial states:} \\
\texttt{```} \\
\texttt{Player 1:} \\
\texttt{ 1234567} \\
\texttt{╔═╤══╤══╗} \\
\texttt{║◌│◌◌│◌◌║ 1} \\
\texttt{║◌└┬─┴──╢ 2} \\
\texttt{║◌◌│L◌◌◌║ 3} \\
\texttt{╟─┬┘◌◌┌─╢ 4} \\
\texttt{║◌│◌◌◌│C║ 5} \\
\texttt{╟─┼─┬─┤P║ 6} \\
\texttt{║◌│◌│◌│◌║ 7} \\
\texttt{╚═╧═╧═╧═╝} \\
\\ 
\\ 
\texttt{Player 2:} \\
\texttt{ 1234567} \\
\texttt{╔═╤══╤══╗} \\
\texttt{║◌│◌◌│◌◌║ 1} \\
\texttt{║L└┬─┴──╢ 2} \\
\texttt{║◌C│◌◌◌◌║ 3} \\
\texttt{╟─┬┘◌◌┌─╢ 4} \\
\texttt{║◌│◌◌◌│◌║ 5} \\
\texttt{╟─┼─┬─┤◌║ 6} \\
\texttt{║◌│P│◌│◌║ 7} \\
\texttt{╚═╧═╧═╧═╝} \\
\\ 
\texttt{```} \\
            }
        }
    }
    & & \\ \\

    \theutterance \stepcounter{utterance}  
    & & & \multicolumn{2}{p{0.3\linewidth}}{
        \cellcolor[rgb]{0.9,0.9,0.9}{
            \makecell[{{p{\linewidth}}}]{
                \texttt{\tiny{[GM$|$GM]}}
                \texttt{End states:} \\
\texttt{```} \\
\texttt{Player 1:} \\
\texttt{ 1234567} \\
\texttt{╔═╤══╤══╗} \\
\texttt{║◌│◌◌│◌◌║ 1} \\
\texttt{║◌└┬─┴──╢ 2} \\
\texttt{║◌◌│◌◌◌◌║ 3} \\
\texttt{╟─┬┘◌◌┌─╢ 4} \\
\texttt{║◌│CLP│◌║ 5} \\
\texttt{╟─┼─┬─┤◌║ 6} \\
\texttt{║◌│◌│◌│◌║ 7} \\
\texttt{╚═╧═╧═╧═╝} \\
\\ 
\\ 
\texttt{Player 2:} \\
\texttt{ 1234567} \\
\texttt{╔═╤══╤══╗} \\
\texttt{║◌│◌◌│◌◌║ 1} \\
\texttt{║◌└┬─┴──╢ 2} \\
\texttt{║◌◌│◌◌◌◌║ 3} \\
\texttt{╟─┬┘◌◌┌─╢ 4} \\
\texttt{║◌│CLP│◌║ 5} \\
\texttt{╟─┼─┬─┤◌║ 6} \\
\texttt{║◌│◌│◌│◌║ 7} \\
\texttt{╚═╧═╧═╧═╝} \\
\\ 
\texttt{```} \\
            }
        }
    }
    & & \\ \\

    \theutterance \stepcounter{utterance}  
    & & & \multicolumn{2}{p{0.3\linewidth}}{
        \cellcolor[rgb]{0.9,0.9,0.9}{
            \makecell[{{p{\linewidth}}}]{
                \texttt{\tiny{[GM$|$GM]}}
                \texttt{* success: True} \\
\texttt{* lose: False} \\
\texttt{* aborted: False} \\
\texttt{{-}{-}{-}{-}{-}{-}{-}} \\
\texttt{* Shifts: 2.00} \\
\texttt{* Max Shifts: 4.00} \\
\texttt{* Min Shifts: 2.00} \\
\texttt{* End Distance Sum: 0.00} \\
\texttt{* Init Distance Sum: 12.67} \\
\texttt{* Expected Distance Sum: 12.57} \\
\texttt{* Penalties: 8.00} \\
\texttt{* Max Penalties: 8.00} \\
\texttt{* Rounds: 12.00} \\
\texttt{* Max Rounds: 12.00} \\
\texttt{* Object Count: 3.00} \\
            }
        }
    }
    & & \\ \\

\end{supertabular}
}

\end{document}
